\section{Conclusion}
\label{sec:conclusion}

We implemented hybrid feature-aware matrix factorization with a WARP ranking
loss end to end, from scratch in NumPy/SciPy: point-in-time feature
engineering over demographic and catalogue attributes, a from-scratch WARP
training loop with the stability corrections it needs in practice, a
dedicated cold-start evaluation protocol that withholds entities from
training entirely rather than merely their last interaction, and a
parity-tested serving path -- precomputed top-$N$ for known entities, a
documented request-time fallback for genuinely new ones -- deployed by
Terraform against a local AWS emulator.

The measured outcome has two parts, and we report both. On warm-start
ranking quality, the simpler ID-only baseline wins, consistent with the
sibling BPR project's own honest negative result: fitting a population
everyone has already been seen in favours fewer structural constraints. On
the cold-start slices this project was built to measure, the comparison is
not close: popularity and ID-only MF are both \emph{structurally} incapable
of saying anything about a withheld entity (exactly zero on every metric),
while the hybrid model retrieves a strictly positive, if modest, signal
through shared brand/category/ingredient structure alone.

\subsection{What transfers}

\begin{itemize}
  \item \textbf{A feature store that is the model, not an accessory to it.}
        Three dedicated test files (schema, leakage, parity) treat the
        feature layer with the same rigor the sibling projects reserve for
        their scoring function, because here they are the same thing.
  \item \textbf{A cold-start harness.} Withholding entities from training
        entirely, restricting identity vocabularies to the training
        partition, and reporting three distinct withheld-population slices
        is a reusable pattern for any feature-aware recommender, not
        specific to skincare.
  \item \textbf{A from-scratch WARP implementation} with the gradient
        derivation and the embedding/bias clipping correction documented in
        Section~\ref{sec:methodology}, removing the fragile \texttt{lightfm}
        Cython dependency entirely.
\end{itemize}

\subsection{Future work}

\paragraph{Sharper negative sampling.} WARP's rejection sampling degrades as
training items increasingly separate into clear winners and losers; sampling
negatives with a mild popularity tilt (as the sibling BPR project's own
future-work notes propose) would very likely close part of the gap to the
ID-only baseline on warm-start quality.

\paragraph{A production-scale catalogue.} The honest caveat in
Section~\ref{sec:results} is partly an artifact of the synthetic fixture's
small catalogue. Running the same pipeline against the real Kaggle dataset
(\texttt{--kaggle}, $\sim$8k products and $\sim$1M reviews) would both sharpen
the ingredient-affinity signal and lower the chance baseline's competitiveness.

\paragraph{Confidence intervals.} Repeating the cold-start comparison across
several random withheld-population draws and reporting intervals would
establish how much of the zero-versus-positive gap is a reliable structural
property versus a property of one particular random split.

\paragraph{Re-admitting popularity as a tempered prior.} Section
\ref{sec:features} excludes \texttt{item\_popularity\_bucket} from the
retrieval representation entirely because an unweighted inclusion suppressed
cold items. A middle ground -- including it with a separately, more
aggressively regularized embedding dimension, or only for warm items at
inference time -- might recover some of the warm-start ranking quality lost
to the ID-only baseline without reintroducing the cold-item suppression
this report documents.
